\section{Proof of Theorem \ref{thm:mainFn}}
\label{sec:main}

In this section, we prove Theorem \ref{thm:mainFn} and \ref{thm:mainFnp2}, completely characterizing Fibonacci polynomials over finite fields that are perfect powers. For the remainder of this paper, let $\bF_q$ denote a finite field of characteristic $p$. Also, we will mainly focus on nonconstant polynomials.

\subsection{Factorization over finite fields}

We first find the factorization of $F_n(T)$.


\begin{lemma}
\label{lem:Fpk}
    Over $\bF_q$, we have
    \begin{equation}
    \label{eqn:Fpkodd}
        F_{p^k}(T) = (T^2 + 4)^{\frac{p^k - 1}{2}}
    \end{equation}
    when \(p\) is odd, and
    \begin{equation}
    \label{eqn:Fpkeven}
        F_{2^k}(T) = T^{2^k - 1}
    \end{equation}
    when \(p = 2\).
\end{lemma}
\begin{proof}
    When $k = 0$, this is trivial. Assume $k \ge 1$.
    When \(p\) is odd, by \eqref{eqn:FnBinet2}
    \[
        F_{p^k}(T) = \frac{\alpha(T)^{p^k} - \beta(T)^{p^k}}{\alpha(T) - \beta(T)} = \frac{(\alpha(T) - \beta(T))^{p^k}}{\alpha(T) - \beta(T)} = (\alpha(T) - \beta(T))^{p^{k} -1} = (T^2 + 4)^{\frac{p^k -1}{2}}.
    \]
    For \(p = 2\), we still have \((\alpha(T) - \beta(T))^2 = (\alpha(T) + \beta(T))^2 = T^2\) and \(\alpha(T) - \beta(T) = T\), so by \eqref{eqn:FnBinet}
    \[
        F_{2^k}(T) = \frac{\alpha(T)^{2^k} - \beta(T)^{2^k}}{\alpha(T) - \beta(T)} = (\alpha(T) - \beta(T))^{2^k - 1} = T^{2^k - 1}.
    \]
    Note that we are working over characteristic $p$, so $(\alpha - \beta)^{p^k} = \alpha^{p^k} - \beta^{p^k}$ holds for any $k \ge 1$.
\end{proof}

\begin{lemma}
\label{lem:Fpkm}
    Write \(n = p^k m\) where \(m\) is coprime to \(p\). Then \(F_n(T)\) factors over \(\bF_q\) as
    \begin{equation}
        F_{p^k m}(T) = F_{p^k}(T) F_m(T)^{p^k}.
    \end{equation}
\end{lemma}
\begin{proof}
    By \eqref{eqn:FnBinet},
    \begin{align*}
        F_n(T) &= \frac{\alpha(T)^{p^k m} - \beta(T)^{p^k m}}{\alpha(T) - \beta(T)} \\
        &= \frac{(\alpha(T)^{m} - \beta(T)^m)^{p^k}}{\alpha(T) - \beta(T)} \\
        &= \frac{(\alpha(T) - \beta(T))^{p^k}}{\alpha(T) - \beta(T)}\left(\frac{\alpha(T)^m - \beta(T)^m}{\alpha(T) - \beta(T)}\right)^{p^k} \\
        &= F_{p^k}(T) F_m(T)^{p^k}.
    \end{align*}
\end{proof}

By Lemma \ref{lem:Fpk} and \ref{lem:Fpkm}, we get:

\begin{proposition}
\label{prop:Fpkm}
    Write \(n = p^k m\) where \(m\) is coprime to \(p\). Then \(F_n(T)\) factors over \(\bF_q\) as
    \begin{equation}
    \label{eqn:Fpkmfactor}
        F_{p^k m}(T) = (T^2 + 4)^{\frac{p^k - 1}{2}} F_m(T)^{p^k}.
    \end{equation}
    When $p = 2$, the above factorization becomes
    \begin{equation}
        \label{eqn:F2kmfactor}
        F_{2^k m}(T) = T^{2^k - 1} F_m(T)^{2^k}.
    \end{equation}
\end{proposition}

\subsection{Square-free factor}

By Proposition \ref{prop:Fpkm}, it is enough to understand the factorization of $F_m(T)$ for $p \nmid m$.
For odd $p$, the following shows that they are square-free. 
% Next, we prove that
\begin{proposition}
\label{prop:Fmsqfree}
    Assume that \(p\) is odd and \(p \nmid m\). Then \(F_m(T)\) is square-free over \(\bF_q\).
\end{proposition}

We provide two different proofs of Proposition \ref{prop:Fmsqfree}.
The first proof is based on the discriminant formula of Florez--Higuita--Ramirez \cite{florez2019resultant}.

\begin{proof}[First proof of Proposition \ref{prop:Fmsqfree}]
    Recall that a polynomial over a field has no multiple root in an algebraic closure if and only if the distriminant is nonzero.
    The discriminant of \(F_m(T)\) is given by \cite[Theorem 4]{florez2019resultant}
    \[
        \mathrm{Disc}(F_m) = (-1)^{\frac{(m-2)(m-1)}{2}} 2^{m-1} m^{m-3}
    \]
    and it is nonzero in \(\bF_q\) when \(2 \ne p \nmid m\), so \(F_m(T)\) is square-free.
\end{proof}

The second proof directly shows that the zeros \eqref{eqn:fibo_zeros} are all distinct modulo $p$ (i.e. distinct in $\overline{\mathbb{F}}_p$).

\begin{lemma} \label{distinct}
Let $p$ be an odd prime and let $n > 1$ and $p \nmid n$. For $k \in \mathbb{Z}$ with $0 < |k| < n$,
$$\zeta_{2n}^k \not \equiv 1 \pmod p,$$
i.e. $\zeta_{2n}^k - 1$ is nonzero in $\overline{\bF}_p$.
\end{lemma}

\begin{proof}
Let:
$$f(T) := \prod_{\substack{-n < k \leq n \\ k \neq 0, \text{ } k \in \mathbb{Z}}} (T - \zeta_{2n}^k) \text{.}$$

When $k = n$, we have $\zeta_{2n}^n = -1 \not \equiv 1 \pmod p$. Suppose $\zeta_{2n}^k \equiv 1 \pmod p$ for some $k$ such that $0 < |k| < n$, then $f(1)\equiv 0\pmod p$.
Notice since:
    \begin{align*}
        T^{2n}-1 &= \prod_{\substack{-n < k \le n \\ k \in \mathbb{Z}}}(T-\zeta_{2n}^k)
        = (T - 1) f(T) \text{,}
    \end{align*}
we have:
\begin{align*}
    f(T) &= \frac{T^{2n} - 1}{T - 1} = \underbrace{T^{2n - 1} + T^{2n - 2} + \dots + T + 1}_{2n \text{ terms}} \text{.}
\end{align*}
From this we get $f(1) = 2n$ and $f(1) \equiv 0 \pmod{p}$, contradicting $p \nmid n$.
\end{proof}

\begin{proof}[Second proof of Proposition \ref{prop:Fmsqfree}]
    If $m = 1$, then $F_1(T) = 1$ is square-free by definition. Assume $m > 1$.
    Suppose $F_m(T)$ is not square-free, so that it has a multiple zero in $\overline{\bF}_p$.
    By \eqref{eqn:fibo_zeros}, there exist distinct $k, l \in \{1, 2, \dots, m - 1\}$ such that $\zeta_{2m}^{k} + \zeta_{2m}^{-k} = \zeta_{2m}^{l} + \zeta_{2m}^{-l}$ in $\overline{\bF}_p$.
    Then $(\zeta_{2m}^{k} - \zeta_{2m}^{l})(\zeta_{2m}^{k + l} - 1) = 0$, hence $\zeta_{2m}^{k} = \zeta_{2m}^{l}$ or $\zeta_{2m}^{k} = \zeta_{2m}^{-l}$.
    By Lemma \ref{distinct}, we have $k \equiv l \pmod{2m}$ or $k \equiv -l \pmod{2m}$, which are both impossible.
\end{proof}

Next, we show that $F_m(T)$ for $p \nmid m$ is always coprime to \(T^2 + 4\), hence $(T^2 + 4)^{\frac{p^k - 1}{2}}$ and $F_m(T)^{p^k}$ in \eqref{eqn:Fpkmfactor} do not share common factors. Knowing that the recurrence relation of the Fibonacci polynomials is \textit{Fibonacci-type}, we can apply the following lemma.

\begin{lemma}[{\cite[Lemma 19]{florez2019resultant}}]
\label{lem:florez_mod}
    For $n \ge 0$, we have
    \begin{equation}
        W_n \pmod{f^2 + 4g} = \begin{cases}
            n (-g)^{(n-1)/2} & 2 \nmid n \\
            (-1)^{(n+2)/2} (nfg^{(n-2)/2}) / 2 & 2 \mid n
        \end{cases}
        ,
    \end{equation}
    where $f$ and $g$ are the coefficient polynomials in \eqref{eqn:lucas_sequence}.
\end{lemma}

\begin{corollary}\label{cor:alpha}
    Let $p$ be an odd prime. We have
    \begin{equation}
        \label{eqn:Fnalpha}
        F_n(T) \pmod{T^2 + 4} = \begin{cases}
            (-1)^{\frac{n-1}{2}}n & 2 \nmid n \\
            (-1)^{\frac{n}{2} + 1} \frac{nT}{2} & 2 \mid n
        \end{cases}
    \end{equation}
    over a field of characteristic $p$.
\end{corollary}
\begin{proof}
    This is a special case of Lemma \ref{lem:florez_mod}, when $f(T) = T$ and $g(T) = 1$.
    Note that 2 is invertible in a field of odd characteristic.
\end{proof}

\begin{corollary}\label{cor:coprime}
     Suppose $p \nmid m$. Then, $T^2 + 4$ and $F_m(T)$ are coprime in $\mathbb{F}_q[T]$.
\end{corollary}
\begin{proof}
    Assume $p$ is an odd prime.
    Suppose $T^2 + 4$ and $F_m(T)$ in $\mathbb{F}_q[T]$ share a factor. Then, there must exist some common zero of the polynomial, which we will denote as $\alpha \in \bF_{q^2}$.
    Then \(\alpha^2 + 4 = 0\) and \(F_m(\alpha) = 0\).
    However, such \(\alpha\) cannot exist by Corollary \ref{cor:alpha}; if \(m\) is even, then \(F_m(\alpha) = (-1)^{\frac{m}{2}+1} \frac{m}{2} \alpha\) and it is nonzero since \(m \ne 0\) in \(\bF_q\), and the case where \(m\) is odd follows similarly.

    For \(p = 2\), \(F_{n}(0) = F_{n-2}(0)\) for all \(n \ge 2\), so \(F_m(0) = F_1(0) = 1\) for all odd \(m\). Hence \(T^2 + 4 = T^2\) and \(F_m(T)\) are coprime for odd \(m\).
\end{proof}

%source
\begin{proposition}
    \label{prop:Fmp2sq}
    Let \(p = 2\). For $n > 0$, \(F_n(T)\) is square in \(\bF_q[T]\) if and only if \(n\) is odd.
\end{proposition}
\begin{proof}
    This immediately follows from \cite[Lemma 3.5]{kitayama2017irreducibility}; we have \(F_{2k}(T) = T F_k(T)^2\) and \(F_{2k+1}(T) = (F_{k+1}(T) + F_k(T))^2\), which can be proven by \eqref{eqn:FnBinet}.
\end{proof}

Finally, before proceeding to the main theorem, we note the following result.

\begin{theorem}
    \label{thm:Fnnonpower}
    Let $p$ be an odd prime number. If $j > 1$ and $p \mid j$, then there exists no $n$ such that $F_n(T)$ is a perfect $j$-th power in $\mathbb{F}_q[T]$.
\end{theorem}

\begin{proof}
    Suppose $j > 1$ and $p \mid j$. Assume there exists some $n = p^km$ with $k \geq 0$ and $p \nmid m$ such that $F_n(T)$ is a perfect $j$-th power in $\mathbb{F}_q[T]$. 
    If $k = 0$, then $p \nmid m$ and $F_n(T) = F_m(T)$ is square-free by Proposition \ref{prop:Fmsqfree}, which cannot be a perfect $j$-th power for $j > 1$, so $k \ge 1$.
    By Lemma \ref{lem:Fpkm}, we have $F_{p^k m}(T) = (T^2 + 4)^{\frac{p^k - 1}{2}} F_m(T)^{p^k}$, where $F_m(T)$ is coprime to $T^2 + 4$ by Corollary \ref{cor:coprime}.
    Then $j$ should divide the exponent $\frac{p^k - 1}{2}$, which contradicts to $p \mid j$.
\end{proof}

\subsection{Proof of the main theorem}

Now we are ready to prove Theorem \ref{thm:mainFn}.

\begin{proof}[Proof of Theorem \ref{thm:mainFn}]
    $(\Leftarrow)$ Let $k \in \mathbb{N}$ be arbitrary. By Lemma \ref{lem:Fpk}, $F_{p^{ak}}(T) = (T^2 +4)^{\frac{p^{ak} - 1}{2}}$, and by our assumption, $p^a \equiv 1 \pmod {2j}$ and thus $\frac{p^{ak} - 1}{2} \equiv 0 \pmod j$. Therefore, $F_n(T)$ where $n = p^{ak}$ is a perfect $j$-th power.
    
    $(\Rightarrow)$ Suppose $F_n(T)$ is a perfect $j$-th power in $\mathbb{F}_q[T]$ and write $n = p^{ak + b}m$, where $p$ and $m$ are coprime (so $p \nmid m$), $k \geq 0$, and \(0 \le b < a\). By Proposition \ref{prop:Fpkm}, we have
    \[
        F_{p^{ak + b}m}(T) = (T^2 + 4)^\frac{p^{ak + b} - 1}{2} \left(F_m(T)\right)^{p^{ak + b}} \text{.}
    \]
    Suppose $b > 0$. Then, $p^{ak + b} = p^{ak}p^b \equiv p^b \pmod{2j}$, where $p^b \not \equiv 1 \pmod{2j}$ by the assumption that $0 < b < a$. This implies that $\frac{p^{ak + b} - 1}{2} \not \equiv 0 \pmod j$, so $(T^2 + 4)^\frac{p^{ak + b} - 1}{2}$ is not a perfect $j$-th power.
    Since \(T^2 + 4\) and \(F_m(T)\) are coprime by Corollary \ref{cor:coprime}, we get a contradiction and we should have \(b = 0\).
    We can rewrite \(F_n(T)\) as
    \[
        F_{p^{ak}m}(T) = (T^2 + 4)^\frac{p^{ak} - 1}{2} \left(F_m(T)\right)^{p^{ak}} \text{,}
    \]
    where \((T^2 + 4)^{\frac{p^{ak} - 1}{2}}\) is a perfect \(j\)-th power and \(F_m(T)\) is square-free by Proposition \ref{prop:Fmsqfree}. If \(m \ne 1\), then \(F_m(T) \ne 1\) and by Corollary \ref{cor:coprime}, the exponent \(p^{ak}\) needs to be a multiple of \(j > 1\), which contradicts to \( \gcd(p, 2j) = 1\).
    This implies \(m = 1\) and we have \(n = p^{ak}\).
\end{proof}

When \(p = 2\), we already showed that \(F_n(T)\) for $n > 0$ is square in \(\bF_q[T]\) if and only if \(n\) is odd. We can further determine higher powers in \(\bF_q[T]\).
\begin{theorem}
\label{thm:mainFnp2}
    Let \(j > 2\) be an odd number and \(a = \mathrm{ord}_j(2)\) be the order of \(2\) in \((\bZ/j\bZ)^\times\).
    Let $q$ be a power of $2$.
    Then \(F_n(T)\) is a perfect \(j\)-th power in $\bF_q[T]$ if and only if \(n = 2^{ak}\) for some \(k \in \bN\).
\end{theorem}
\begin{proof}
    If $n = 2^{ak}$, then $2^a \equiv 1 \pmod{j}$ and \eqref{eqn:Fpkeven} shows that $F_{n}(T) = T^{2^{ak} - 1} = (T^{\frac{2^{ak} - 1}{j}})^{j}$ is a perfect $j$-th power.
    Conversely, assume that $n = 2^{s} \cdot m$ for $s \ge 0$ and odd $m$.
    From \eqref{eqn:fibo_poly}, one can easily check that $F_m(0) = 1$ for all odd $m$.
    In particular, $\gcd(T, F_m(T)) = 1$.
    If $m > 1$, then by \eqref{eqn:F2kmfactor} $F_n(T)$ cannot be a perfect $j$-th power.
    Hence $m = 1$ and $F_n(T) = F_{2^s}(T) = T^{2^{s} - 1}$, which is a perfect $j$-th power if and only if $2^s \equiv 1 \pmod{j} \Leftrightarrow s \equiv 0 \pmod{a}$.
\end{proof}

Similarly, we can also prove Theorem \ref{thm:mainFnpowerful}.
\begin{proof}[Proof of Theorem \ref{thm:mainFnpowerful}]
    $(\Rightarrow)$ If $p \nmid n$, then $F_n(T)$ is square-free and cannot be a powerful polynomial.

    $(\Leftarrow)$ If $n = p^k m$ with $k \ge 1$, then $F_n(T)$ is powerful by \eqref{eqn:Fpkmfactor}.
    Here we use $p > 3$ to guarantee $\frac{p^k - 1}{2} > 1$.
\end{proof}

We need extra care for $p = 3$, since we can have $\frac{p^k - 1}{2} = 1$ when $k = 1$.
\begin{corollary}
    \label{cor:powerfulp3}
    For $p = 3$, $F_n(T)$ is powerful over $\bF_q$ if and only if $9 \mid n$.
\end{corollary}
\begin{proof}
    The proof is similar to that of Theorem \ref{thm:mainFnpowerful}, except that we need $k \ge 2$ to have $\frac{3^k - 1}{2} > 1$.
\end{proof}

Similar characterization is also available for $p = 2$.
\begin{corollary}
\label{cor:powerfulp2}
    For $p = 2$, $F_n(T)$ is powerful over $\bF_q$ if and only if $n \ge 3$ is odd or a multiple of $4$.
\end{corollary}
\begin{proof}
    By \eqref{eqn:F2kmfactor} and \cite[Lemma 3.5]{kitayama2017irreducibility}, when $n = 2^k m$ with odd $m$, we have 
    \[F_n(T) = T^{2^k - 1}F_{m}(T)^{2^k} = T^{2^k - 1}(F_{(m+1)/2}(T)  + F_{(m-1)/2}(T))^{2^{k+1}}.\]
    If $k = 0$ and $n = m \ge 3$ is odd, then $F_n(T) = (F_{(m+1)/2}(T)  + F_{(m-1)/2}(T))^{2}$ is a square of a non-constant polynomial, hence powerful.
    Similarly, $F_n(T)$ is powerful when $k \ge 2$.
    However, if $k = 1$, then $F_n(T) = TF_m(T)^2$ is not powerful, since $T$ and $F_m(T)$ are coprime and $F_n(T)$ is divisible by $T$ only once.
\end{proof}

Table \ref{tab:fibo} shows the factorizations of $F_n(T)$ over $\bF_p$ for $p = 2, 3, 5$ and $1 \le n \le 20$, which is consistent with the above characterizations.